\subsection{Especificaciones de los módulos críticos}

\subsubsection*{M1: Captura e ingesta}

La Tabla~\ref{tab:modulo_m1} especifica el punto de entrada de eventos clínicos.

\begin{table}[H]
    \caption{Especificación arquitectónica de M1}
    \label{tab:modulo_m1}
    \begin{tabularx}{\textwidth}{|l|X|}
        \hline
        \textbf{Atributo} & \textbf{Descripción} \\
        \hline
        Propósito & Recibir texto o audio desde Flutter e importar recursos clínicos sintéticos desde un sistema hospitalario simulado. \\
        \hline
        Flujo objetivo & Validar sesión y caso para la captura asistencial; validar y normalizar recursos FHIR o mensajes HL7v2 para la ingesta simulada; y asignar metadatos antes de persistir cada resultado. \\
        \hline
        Dependencias & Flutter, FastAPI, servicio ASR, adaptadores FHIR/HL7v2, PostgreSQL y Storage. \\
        \hline
        Restricciones & Validar formato y tamaño del audio, rechazar recursos clínicos inválidos, utilizar solo datos sintéticos y no confiar en identificadores enviados sin validación. \\
        \hline
    \end{tabularx}
\end{table}

M1 admite transcripción local o remota sin alterar el contrato del evento. Los audios remotos se conservan en Storage y la identidad del autor se obtiene de la sesión validada por FastAPI. En el flujo de interoperabilidad, los adaptadores transforman datos sintéticos al modelo interno antes de dejarlos disponibles para consulta.

\subsubsection*{M2: Privacidad y desidentificación}

La Tabla~\ref{tab:modulo_m2} define la barrera previa al procesamiento generativo.

\begin{table}[H]
    \caption{Especificación arquitectónica de M2}
    \label{tab:modulo_m2}
    \begin{tabularx}{\textwidth}{|l|X|}
        \hline
        \textbf{Atributo} & \textbf{Descripción} \\
        \hline
        Propósito & Detectar y reemplazar IPS antes de enviar texto a un LLM externo. \\
        \hline
        Flujo objetivo & Aplicar reglas y NER, conservar la relación de sustitución en un entorno protegido y entregar solo texto sanitizado a M3. \\
        \hline
        Dependencias & M1, motor NER y políticas institucionales de tratamiento de datos. \\
        \hline
        Restricciones & Medir precisión y exhaustividad con datos sintéticos antes de procesar información real. \\
        \hline
    \end{tabularx}
\end{table}

La separación de M2 evita que el proveedor generativo reciba identificadores innecesarios. Su procesamiento combina reglas y NER antes de entregar texto sanitizado a M3~\cite{nlp_ner_privacy}.

\subsubsection*{M3: Procesamiento cognitivo}

La Tabla~\ref{tab:modulo_m3} resume las restricciones aplicables a la clasificación y síntesis.

\begin{table}[H]
    \caption{Especificación arquitectónica de M3}
    \label{tab:modulo_m3}
    \begin{tabularx}{\textwidth}{|l|X|}
        \hline
        \textbf{Atributo} & \textbf{Descripción} \\
        \hline
        Propósito & Sugerir categorías y producir borradores de evolución a partir de eventos del turno. \\
        \hline
        Flujo objetivo & Recibir texto sanitizado, aplicar instrucciones acotadas y devolver una propuesta identificada como contenido generado. \\
        \hline
        Dependencias & M2, servicio LLM y catálogo de categorías. \\
        \hline
        Restricciones & No diagnosticar, no completar antecedentes ausentes y no publicar resultados automáticamente. \\
        \hline
    \end{tabularx}
\end{table}

M3 opera como asistencia y no como fuente autónoma de verdad clínica. La literatura evidencia errores residuales en la generación de resúmenes, lo que obliga a mantener revisión experta~\cite{rust2025}.

\subsubsection*{M4: Gobernanza y validación clínica}

La Tabla~\ref{tab:modulo_m4} describe el control HITL del ciclo de vida de cada evento.

\begin{table}[H]
    \caption{Especificación arquitectónica de M4}
    \label{tab:modulo_m4}
    \begin{tabularx}{\textwidth}{|l|X|}
        \hline
        \textbf{Atributo} & \textbf{Descripción} \\
        \hline
        Propósito & Exigir revisión, corrección y confirmación antes de consolidar contenido generado por IA. \\
        \hline
        Flujo objetivo & Recibir borrador, presentarlo como editable, registrar la decisión del profesional y publicar solo tras aprobación. \\
        \hline
        Dependencias & M3, M5, M6 y aplicación Flutter. \\
        \hline
        Restricciones & Impedir autoaprobación y distinguir claramente contenido preliminar de contenido confirmado. \\
        \hline
    \end{tabularx}
\end{table}

El flujo HITL mantiene la responsabilidad en el profesional y permite corregir errores previsibles de la IA antes de integrar un evento a la bitácora~\cite{lazaros2026}.

\subsubsection*{M5 y M6: Identidad y auditoría}

La Tabla~\ref{tab:modulos_m5_m6} reúne las responsabilidades transversales de control de acceso y trazabilidad.

\begin{table}[H]
    \caption{Especificación arquitectónica de M5 y M6}
    \label{tab:modulos_m5_m6}
    \begin{tabularx}{\textwidth}{|l|X|X|}
        \hline
        \textbf{Módulo} & \textbf{Propósito} & \textbf{Condición objetivo} \\
        \hline
        M5: Identidad & Autenticar, validar sesiones y aplicar permisos por rol e institución. & Supabase Auth emite la sesión y FastAPI valida cada operación protegida. \\
        \hline
        M6: Auditoría & Registrar actor, acción, recurso, resultado y marca temporal. & Los registros se protegen frente a cambios y se someten a una política de retención. \\
        \hline
    \end{tabularx}
\end{table}

Supabase Auth administra las sesiones y FastAPI valida el usuario y sus permisos antes de procesar operaciones clínicas. La auditoría registra las acciones críticas junto con el actor, el recurso, el resultado y la marca temporal.
